% Source printed pp. 426-430 / PDF pp. 434-438.
\section{GIT 的例子}\label{src:8.5}
\begin{epigraph}
要透彻理解任何概念, 都离不开尽可能丰富的好例子. 我想学习新事物时, 总把积累这样一批例子作为首要工作.
\epigraphauthor{Paul Halmos}
\end{epigraph}
仿射 GIT 的例子见例~\bookref{7.5.7}--\bookref{7.5.8}, 射影 GIT 的初等例子见例~\bookref{8.2.8}--\bookref{8.2.7}.
\subsection{射影空间中的有序点}\label{src:8.5.1-subsection}
\begin{example}\label{src:8.5.1}
考虑 $\mathrm{SL}_2$ 在 $X=(\mathbb P^1)^n$ 上的对角作用, 以及由带 $\mathrm{SL}_2$-作用的线丛 $\mathcal O(1)\boxtimes\cdots\boxtimes\mathcal O(1)$ 定义的 $\mathrm{SL}_2$-等变 Segre 嵌入 $(\mathbb P^1)^n\to\mathbb P^{2^n-1}$. 我们断言
\begin{align*}
X^{\mathrm s}&=\{(p_1,\ldots,p_n)\mid\text{对所有 }q\in\mathbb P^1,\ \#\{i\mid p_i=q\}<n/2\},\\
X^{\mathrm{ss}}&=\{(p_1,\ldots,p_n)\mid\text{对所有 }q\in\mathbb P^1,\ \#\{i\mid p_i=q\}\leq n/2\}.
\end{align*}
为证明这一点, 设 $(p_1,\ldots,p_n)\in X(\Bbbk)$, $\lambda:\mathbb G_m\to\mathrm{SL}_2$ 是单参数子群. 存在 $g\in\mathrm{SL}_2(\Bbbk)$, 使对某个 $d\in\mathbb Z$ 有 $g\lambda g^{-1}=\lambda_0^d$, 其中 $\lambda_0(t)=\begin{pmatrix}t^{-1}&0\\0&t\end{pmatrix}$. 可设 $d\geq0$, 因为 $d<0$ 可类似处理. 由 $\mu(x,\lambda)=\mu(gx,\lambda_0^d)=d\mu(gx,\lambda_0)$, 只须计算 $\mu(gx,\lambda_0)$. $\mu(-,\lambda_0)$ 关于 $S_n$-作用对称, 故可设 $gx=(0,\ldots,0,p_k,\ldots,p_n)$, $p_k,\ldots,p_n\ne0$. Segre 嵌入的坐标形如 $(\prod_{i\in\Sigma}x_i)(\prod_{i\notin\Sigma}y_i)$, 其中 $\Sigma\subseteq\{1,\ldots,n\}$, 其权为 $n-2(\#\Sigma)$. 在 $gx$ 处非零且权最大的坐标是 $y_1\cdots y_kx_{k+1}\cdots x_n$, 权为 $2k-n$. 所以 $\mu(gx,\lambda_0)=n-2k$. 因此, 若至多 (分别地, 少于) $n/2$ 个点 $p_i$ 相同, 则 $x$ 半稳定 (分别地, 稳定) 当且仅当 $n\geq2k$ (分别地, $n>2k$). 反之, 若多于 (分别地, 至少) $n/2$ 个点相同, 则经 $\mathrm{SL}_2$ 的一个元素平移并利用 $S_n$-作用的对称性, 可写 $u=(0,\ldots,0,p_k,\ldots,p_n)$, $k>n/2$ (分别地, $k\geq n/2$), 而 $\lambda_0=\operatorname{diag}(t^{-1},t)$ 使 $u$ 不稳定化.

若 $n$ 为奇数, 则 $X^{\mathrm{ss}}=X^{\mathrm s}$, 且 $X^{\mathrm{ss}}\to X^{\mathrm{ss}}\sslash\mathrm{SL}_2$ 是几何商. 若 $n$ 为偶数, 映射 $X^{\mathrm{ss}}\to X^{\mathrm{ss}}\sslash\mathrm{SL}_2$ 将 $(p_1,\ldots,p_n)$ 与 $(q_1,\ldots,q_n)$ 等同, 只要存在大小为 $n/2$ 的子集 $\Sigma\subseteq\{1,\ldots,n\}$, 使对所有 $i,j\in\Sigma$ 有 $p_i=p_j$ 且 $q_i=q_j$. 此时纤维中的唯一闭轨道是这样一个 $n$ 元组的轨道: $\Sigma$ 中位置取 $0$, 其他位置取 $\infty$. 补集 $X^{\mathrm{ss}}\setminus X^{\mathrm s}$ 恰有 $\frac12\binom n{n/2}$ 个闭轨道. 因为 $\operatorname{Aut}(\mathbb P^1)=\mathrm{PGL}_2$, 商 $X^{\mathrm{ss}}\sslash\mathrm{SL}_2=X^{\mathrm{ss}}\sslash\mathrm{PGL}_2$ 可视为 $\mathbb P^1$ 中 $n$ 个有序点在射影等价意义下的模空间的紧化.
\end{example}
\begin{exercise}\label{src:8.5.2}
\begin{enumerate}[(a)]
\item 对 $\mathrm{SL}_2$ 在次数 $n$ 的二元形式的射影化 $\mathbb P(\Gamma(\mathbb P^1,\mathcal O(n)))\cong\mathbb P^n$ 上的作用, 证明半稳定 (分别地, 稳定) 轨迹由每个一次因子的重数小于等于 (分别地, 小于) $n/2$ 的二元形式 $f(x,y)$ 组成.
\item 对 $(\mathbb P^1)^n$ 上的 $\mathrm{SL}_2$-线性化 $\mathcal O(a_1)\boxtimes\cdots\boxtimes\mathcal O(a_n)$, 其中每个 $a_i>0$, 证明半稳定 (分别地, 稳定) 轨迹由满足下列条件的元组 $(p_1,\ldots,p_n)$ 组成: 对所有 $q\in\mathbb P^1(\Bbbk)$,
\[
\sum_{p_i=q}a_i\leq\frac12\sum_{i=1}^na_i\qquad(\text{分别地, }<).
\]
\item 对 $\mathrm{SL}_{r+1}$ 在 $(\mathbb P^r)^n$ 上的作用和 $\mathrm{SL}_{r+1}$-线性化 $\mathcal O(a_1)\boxtimes\cdots\boxtimes\mathcal O(a_n)$, 其中每个 $a_i>0$, 证明半稳定 (分别地, 稳定) 轨迹由满足下列条件的元组 $(p_1,\ldots,p_n)$ 组成: 对每个线性子空间 $W\subsetneq\mathbb P^r$,
\[
\sum_{p_i\in W}a_i\leq\frac{\dim W+1}{r+1}\left(\sum_{i=1}^na_i\right)\qquad(\text{分别地, }<).
\]
\end{enumerate}
\end{exercise}
\subsection{超曲面的稳定性}\label{src:8.5.2-subsection}
\begin{exercise}[三次曲线]\label{src:8.5.3}
考虑 $\mathrm{SL}_3$ 在 $\mathbb P^2$ 中三次曲线的射影空间 $\mathbb P(\mathrm H^0(\mathbb P^2,\mathcal O(3)))$ 上的作用. 证明稳定 (分别地, 半稳定) 轨迹由光滑 (分别地, 至多具有结点的) 曲线组成. 另见 \cite[\S1.11]{Mum77}, \cite[\S4.4]{New78} 和 \cite[第 202--206 页]{HM98}.
\end{exercise}
\begin{exercise}[四次曲线, 困难]\label{src:8.5.4}
更复杂的计算表明, 对 $\mathrm{SL}_3$ 在 $\mathbb P(\mathrm H^0(\mathbb P^2,\mathcal O(4)))$ 上的作用, 四次曲线半稳定当且仅当它不含三重点, 且不是三次曲线与拐点切线之并; 它稳定当且仅当至多具有结点和尖点奇点. 另见 \cite[\S1.12--13]{Mum77} 和 \cite[\S4A]{HM98}.
\end{exercise}
\begin{remark}[三次曲面]\label{src:8.5.5}
对 $\mathrm{SL}_4$ 在 $\mathbb P(\mathrm H^0(\mathbb P^3,\mathcal O(3)))$ 上的作用, 三次曲面稳定 (分别地, 半稳定) 当且仅当它的奇点有限, 且奇点均为普通二重点 (分别地, 普通二重点或轴不包含在曲面中的秩二二重点). 见 \cite[定理 7.14]{Muk03}.
\end{remark}
\subsection{箭图与环面 GIT}\label{src:8.5.3-subsection}
\begin{exercise}[箭图 GIT]\label{src:8.5.6}
\newterm{箭图}{quiver} $Q=(Q_0,Q_1)$ 是有向图, 其中 $Q_0$ 是有限顶点集, $Q_1$ 是有限箭头集; 有源与靶映射 $s,t:Q_1\to Q_0$. $Q$ 的\newterm{箭图表示}{quiver representation}由每个 $i\in Q_0$ 的向量空间 $V_i$ 及每个箭头 $\alpha:i\to j$ 的线性映射 $L_\alpha:V_i\to V_j$ 组成. 若每个 $V_i$ 有限维, $d_i=\dim V_i$, 则称 $d=(d_i)$ 为 $V$ 的\newterm{维数向量}{dimension vector}.

固定维数向量 $d=(d_i)\in\mathbb Z^{Q_0}$, 考虑维数向量为 $d$ 的表示空间
\[
R(Q,d)=\prod_{\alpha\in Q_1}\Hom(\Bbbk^{d_{s(\alpha)}},\Bbbk^{d_{t(\alpha)}}).
\]
它通过 $(g_i)\cdot(L_\alpha)=(g_{t(\alpha)}L_\alpha g_{s(\alpha)}^{-1})$ 继承 $\prod_i\mathrm{GL}_{d_i}$ 的作用. 由标量矩阵元组 $(t\,\mathrm{id}_{\Bbbk^{d_i}})$ ($t\in\mathbb G_m$) 组成的对角子群 $\mathbb G_m\subseteq\prod_i\mathrm{GL}_{d_i}$ 正规且作用平凡, 故商 $G:=(\prod_i\mathrm{GL}_{d_i})/\mathbb G_m$ 也作用于 $R(Q,d)$.

对满足 $\sum_i a_id_i=0$ 的任意整数元组 $a=(a_i)_{i\in Q_0}$, 考虑特征标
\[
\chi_a:G\longrightarrow\mathbb G_m,\qquad(g_i)\longmapsto\prod_i\det(g_i)^{a_i}.
\]
用仿射 Hilbert--Mumford 判据 (\bookref{8.4.9}) 证明, 表示 $V\in R(Q,d)$ 关于 $\chi$ 半稳定 (分别地, 稳定), 当且仅当对每个子表示 $W\subseteq V$ (即子空间 $W_i\subseteq V_i$ 满足 $L_\alpha(W_{s(\alpha)})\subseteq W_{t(\alpha)}$), 有
\[
\sum_i a_i\dim W_i\geq0\qquad(\text{分别地, }>).
\]
另见 \cite[命题 3.1]{Kin94}.
\end{exercise}
\begin{remark}[环面簇的 Cox 构造]\label{src:8.5.7}
设 $X=X(\Sigma)$ 是具有扇 $\Sigma\subseteq N_{\mathbb R}$ 和环面 $T_N$ 的固有环面簇, 其中 $N$ 是格, 对偶为 $M$. 以 $\Sigma(1)$ 表示扇的射线, 则与 $\rho\in\Sigma(1)$ 对应的除子 $D_\rho$ 生成类群. 有短正合列
\[
0\longrightarrow M\longrightarrow\mathbb Z^{\Sigma(1)}\longrightarrow\operatorname{Cl}(X)\longrightarrow0.
\]
代数群 $G:=\Hom(\operatorname{Cl}(X),\mathbb G_m)$ 可对角化 (因而线性约化), 且属于对上列施加 $\Hom(-,\mathbb G_m)$ 所得的短正合列
\[
1\longrightarrow G\longrightarrow\mathbb G_m^{\Sigma(1)}\longrightarrow T_N\longrightarrow1.
\]
群 $G$ 自然作用于 $\mathbb A^{\Sigma(1)}$.

对锥 $\sigma\in\Sigma$, 令 $x^\sigma:=\prod_{\rho\notin\sigma(1)}x_\rho$. 以 $\sigma$ 遍历最大维锥时单项式 $x^\sigma$ 生成的理想之零点定义闭 $G$-不变子概形 $Z\subseteq\mathbb A^{\Sigma(1)}$. 此集合也可描述为并 $\bigcup_CV(x_\rho\mid\rho\in C)$, 其中 $C$ 遍历本原集合 $C\subseteq\Sigma(1)$, 即不包含于任何 $\sigma\in\Sigma$ 的 $\sigma(1)$ 中, 而对每个 $C'\subsetneq C$ 都存在 $\sigma\in\Sigma$ 使 $C'\subseteq\sigma(1)$ 的子集.

这里的主要定理是 $X$ 同构于良商 $(\mathbb A^{\Sigma(1)}\setminus Z)\sslash G$. 这就是 $X$ 的 \textit{Cox 构造}, 它为 $X$ 提供齐次坐标, 类似于 $\mathbb A^{n+1}$ 为 $\mathbb P^n=(\mathbb A^{n+1}\setminus0)/\mathbb G_m$ 提供齐次坐标. 当 $\Sigma$ 是单纯扇时, $X$ 是几何商 $(\mathbb A^{\Sigma(1)}\setminus Z)/G$. 此外, 类群 $\operatorname{Cl}(X)$ 与特征标群 $\mathbb X^*(G)$ 等同; 若 $L$ 是 $X$ 上对应于特征标 $\chi$ 的丰沛线丛, 则 $\mathbb A^{\Sigma(1)}\setminus Z$ 是 $G$ 在 $\mathbb A^{\Sigma(1)}$ 上关于 $\chi$ 的半稳定轨迹. 见 \cite{Cox95} 和 \cite[\S5]{CLS11}.
\end{remark}
\subsection{GIT 与双有理几何}\label{src:8.5.4-subsection}
\begin{example}[\texorpdfstring{$\mathbb G_m$}{Gm}-作用的 GIT 变分]\label{src:8.5.8}
考虑 $\mathbb G_m$ 在 $\Bbbk$ 上有限型仿射概形 $X=\Spec A$ 上的作用. 在此例中, 我们研究改变 $\mathbb G_m$ 的特征标时 GIT 商如何变化; 关于特征标的仿射 GIT 见习题~\bookref{8.2.10}. 有双射 $\Hom(\mathbb G_m,\mathbb G_m)\cong\mathbb Z$, 以 $\chi_d(t)=t^d$ 表示与 $d\in\mathbb Z$ 对应的特征标. 设 $A=\bigoplus_{n\in\mathbb Z}A_n$ 是诱导分次. 关于 $\chi_d$ 的半稳定轨迹 $X_{\chi_d}^{\mathrm{ss}}$ 有三种情形:
\begin{enumerate}[(i)]
\item $d=0$: $X^{\mathrm{ss}}(0):=X_{\chi_0}^{\mathrm{ss}}=X$, 且 $X_{\chi_0}^{\mathrm{ss}}\sslash\mathbb G_m=\Spec A_0$.
\item $d>0$: $X^{\mathrm{ss}}(+):=X_{\chi_d}^{\mathrm{ss}}=X\setminus V(\sum_{n>0}A_n)$, 而 $X_{\chi_d}^{\mathrm{ss}}=\Proj\bigoplus_{n\geq0}A_{nd}$ 与 $d$ 无关; 此外, $X^{\mathrm{ss}}(0)$ 与单参数子群 $\chi_d$ 所给出的 $X_{\chi_d}^+$ 等同 (习题~\bookref{7.8.7}).
\item $d<0$: $X^{\mathrm{ss}}(-):=X_{\chi_d}^{\mathrm{ss}}=X\setminus V(\sum_{n<0}A_n)=X_{\chi_d}^+$, 且 $X_{\chi_d}^{\mathrm{ss}}=\Proj\bigoplus_{n\geq0}A_{-nd}$ 与 $d$ 无关.
\end{enumerate}
有交换图
\[
\begin{tikzcd}[ampersand replacement=\&]
X^{\mathrm{ss}}(+)\arrow[r,hook]\arrow[d]\&X\arrow[d]\&X^{\mathrm{ss}}(-)\arrow[l,hook']\arrow[d]\\
{X^{\mathrm{ss}}(+)\sslash\mathbb G_m}\arrow[r]\&{X\sslash\mathbb G_m}\&{X^{\mathrm{ss}}(-)\sslash\mathbb G_m}\arrow[l]
\end{tikzcd}
\]
其中竖直映射是良商, 上方映射是开浸入, 下方映射射影. 仿射 Hilbert--Mumford 判据 (\bookref{8.4.9}) 给出关于 $\chi_0$, $\chi_1$, $\chi_{-1}$ 的稳定轨迹的等同:
\begin{align*}
X^{\mathrm s}(0)&=X\setminus(X^{\mathrm{ss}}(+)\cap X^{\mathrm{ss}}(-)),\\
X^{\mathrm s}(+)&=X^{\mathrm{ss}}(+)=X\setminus X^{\mathrm{ss}}(-),\\
X^{\mathrm s}(-)&=X^{\mathrm{ss}}(-)=X\setminus X^{\mathrm{ss}}(+).
\end{align*}
因此, 若 $X^{\mathrm s}(+)$ 与 $X^{\mathrm s}(-)$ 均非空, 则 $X^{\mathrm{ss}}(+)\sslash\mathbb G_m\to X\sslash\mathbb G_m$ 和 $X^{\mathrm{ss}}(-)\sslash\mathbb G_m\to X\sslash\mathbb G_m$ 在 $X^{\mathrm s}(0)/\mathbb G_m$ 上是同构, 特别地, 它们双有理. 又若 $X^{\mathrm{ss}}(+)$ 与 $X^{\mathrm{ss}}(-)$ 在 $X$ 中的补集余维均至少为二, 则双有理映射 $X^{\mathrm{ss}}(+)\sslash\mathbb G_m\dashrightarrow X^{\mathrm{ss}}(-)\sslash\mathbb G_m$ 在余维 $2$ 意义下是同构, 且除子 $\mathcal O(1)$ (它在 $X\sslash\mathbb G_m$ 上相对丰沛) 推前为 $X^{\mathrm{ss}}(-)\sslash\mathbb G_m$ 上的除子, 其对偶相对丰沛; 即该双有理映射是关于 $\mathcal O(1)$ 的\newterm{翻转}{flip}.
\end{example}
\begin{remark}[GIT 变分]\label{src:8.5.9}
上例可推广到线性约化群 $G$ 的作用. 推广代数无关的概念, 若存在连通簇 $T$, 点 $t_1,t_2\in T(\Bbbk)$ 以及 $X\times T$ 上带 $G$-作用的线丛 $\mathcal L$, 使 $L_i\cong\mathcal L|_{X\times\{t_i\}}$, 则称射影概形 $X$ 上带 $G$-作用的两个线丛 $L_1,L_2$ 为 $G$-代数等价. $X$ 上带 $G$-作用的线丛模 $G$-代数等价所成的 $G$-等变 Néron--Severi 群 $\operatorname{NS}^G(X)$ 有限生成. $\operatorname{NS}^G(X)_{\mathbb R}\to\operatorname{NS}(X)$ 的核与有理特征标群 $\mathbb X^*(G)_{\mathbb R}$ 等同. 以 $\operatorname{Eff}^G(X)\subseteq\operatorname{NS}^G(X)_{\mathbb R}$ 表示 $G$-有效线性化所成的锥, 即对某个 $d>0$ 存在 $L^{\otimes d}$ 的非零不变截面的 $G$-线性化 $L$, 换言之, $X_L^{\mathrm{ss}}\ne\emptyset$. 又以 $\operatorname{Amp}^G(X)\subseteq\operatorname{NS}^G(X)_{\mathbb R}$ 表示丰沛 $G$-线性化的锥.

GIT 变分的主要结果可表述如下. 半稳定轨迹 $X_L^{\mathrm{ss}}$ 只依赖于 $L$ 的 $G$-代数等价类. 锥 $\operatorname{Amp}^G(X)\cap\operatorname{Eff}^G(X)$ 有由余维 $1$ 的墙确定的多面体分解, 使半稳定轨迹在每个开室内不变. 若 $L_0$ 位于墙上, 而 $L_+$ 与 $L_-$ 位于两侧相邻的室内, 则有交换图
\[
\begin{tikzcd}[ampersand replacement=\&]
X_{L_+}^{\mathrm{ss}}\arrow[r,hook]\arrow[d]\&X_{L_0}^{\mathrm{ss}}\arrow[d]\&X_{L_-}^{\mathrm{ss}}\arrow[l,hook']\arrow[d]\\
{X_{L_+}^{\mathrm{ss}}\sslash G}\arrow[r]\&{X_{L_0}^{\mathrm{ss}}\sslash G}\&{X_{L_-}^{\mathrm{ss}}\sslash G}\arrow[l]
\end{tikzcd}
\]
其中竖直映射是良商, 上方映射是开浸入, 下方映射射影. 若 $X_{L_+}^{\mathrm s}$ 与 $X_{L_-}^{\mathrm s}$ 非空, 则下方映射双有理; 当这些映射在余维 $2$ 意义下是同构时, $X_{L_+}^{\mathrm{ss}}\sslash G\dashrightarrow X_{L_-}^{\mathrm{ss}}\sslash G$ 是关于 $X_{L_+}^{\mathrm{ss}}\sslash G$ 上线丛 $\mathcal O(1)$ 的翻转, 该线丛在 $X_{L_0}^{\mathrm{ss}}\sslash G$ 上相对丰沛. 见 \cite{Tha96,DH98}.
\end{remark}
\begin{remark}[Mori 梦空间]\label{src:8.5.10}
Mori 纲领与 GIT 变分有一种有趣的联系. 正规 $\mathbb Q$-阶乘射影簇 $X$ 称为\newterm{Mori 梦空间}{Mori dream space}, 若:
\begin{enumerate}[(1)]
\item $\Pic(X)_{\mathbb Q}=\operatorname{NS}(X)_{\mathbb Q}$;
\item 数值有效线丛的锥 $\operatorname{Nef}(X)$ 是有限个半丰沛线丛的仿射包;
\item 存在有限个到 $\mathbb Q$-阶乘正规射影簇 $X_i$ 的双有理映射 $f_i:X\dashrightarrow X_i$, 它们在余维 $1$ 意义下是同构, 使可动锥 $\operatorname{Mov}(X)$ 是 $f_i^{-1}(\operatorname{Amp}(X_i)_{\mathbb Q})$ 之并 (若线丛的稳定基点轨迹的余维至少为 $2$, 则称之为\newterm{可动}{movable}).
\end{enumerate}
换言之, 若 $\operatorname{Mov}(X)$ 有有限的墙与室分解, 使线丛确定的射影簇在每个开室内不变, 则 $X$ 是 Mori 梦空间. 等价地, $X$ 是 Mori 梦空间当且仅当 $\Pic(X)_{\mathbb Q}=\operatorname{NS}(X)_{\mathbb Q}$ 且 \textit{Cox 环}
\[
\operatorname{Cox}(X):=\bigoplus_{(d_1,\ldots,d_n)\in\mathbb N^n}\Gamma(X,L_1^{d_1}\otimes\cdots\otimes L_n^{d_n})
\]
有限生成, 其中 $L_1,\ldots,L_n$ 是 $\Pic(X)_{\mathbb Q}$ 的一组基, 其仿射包包含 $\operatorname{Eff}(X)_{\mathbb Q}$. 若 $X$ 是 Mori 梦空间, 则每个双有理模型 $X_i$ 都是 $\Spec(\operatorname{Cox}(X))$ 关于某特征标的半稳定轨迹对环面 $\mathbb G_m^n$ 的 GIT 商. 此外, $\operatorname{Mov}(X)$ 的 Mori 室与 $\mathbb G_m^n$ 在 $\Spec(\operatorname{Cox}(X))$ 上作用的 GIT 变分室相对应. 见 \cite{HK00}.
\end{remark}
\begin{example}[部分消奇]\label{src:8.5.11}
若 $U$ 是光滑簇, $U\to X$ 是对线性约化群的几何商, 则 $X$ 必具有有限商奇点; 这由局部结构定理 (\bookref{5.4.13}) 得到. 另一方面, 若 $U\to X$ 是良商, 则 $X$ 的奇点可能更坏 (\S\bookref{7.5.11-subsection}). 不过, 若存在开子集 $X'\subseteq X$, 使 $\pi_0(X')\to X'$ 是几何商 (例如, $U=V^{\mathrm{ss}}$ 是 $G$-等变嵌入 $V\hookrightarrow\mathbb P^n$ 的半稳定轨迹, 且稳定轨迹 $V^{\mathrm s}$ 非空), 则有典范过程部分消解 $X$ 的奇点, 使其成为有限商奇点. 具体地, 存在交换图
\[
\begin{tikzcd}[ampersand replacement=\&]
U_n\arrow[r]\arrow[d,"\pi_n"']\&U_{n-1}\arrow[r]\arrow[d,"\pi_{n-1}"']\&\cdots\arrow[r]\&{U=U_0}\arrow[d,"\pi_0"]\\
X_n\arrow[r]\&X_{n-1}\arrow[r]\&\cdots\arrow[r]\&{X=X_0}
\end{tikzcd}
\]
使得:
\begin{itemize}
\item 每个 $U_{i+1}$ 都是爆破 $\operatorname{Bl}_{Z_i}U_i$ 的 $G$-不变开子概形, 其中 $Z_i$ 是稳定子维数最大的 $G$-不变光滑闭子概形, $U_{i+1}\subseteq\operatorname{Bl}_{Z_i}U_i$ 是 $\pi_i^{-1}(\pi_i(Z_i))$ 的严格变换之补集. 若 $U=V^{\mathrm{ss}}$ 是射影簇关于 $G$-线性化 $L$ 的半稳定轨迹, 则 $U_{i+1}$ 是关于 $(q^*L)^{\otimes n}\otimes\mathcal O(-E)$ ($n\gg0$) 的半稳定轨迹, 其中 $q:\operatorname{Bl}_{Z_i}U_i\to U_i$, $E$ 为例外除子.
\item 映射 $X_{i+1}\to X_i$ 射影且双有理.
\item 映射 $\pi_i:U_i\to X_i$ 是对 $G$ 的良商, $\pi_n:U_n\to X_n$ 是几何商. 特别地, $X_n$ 具有有限商奇点.
\end{itemize}
见 \cite{Kir85,Rei89,ER21}.

作为此过程的简单例子, 考虑权为 $1$ 和 $-1$ 的 $\mathbb G_m$ 在 $\mathbb A^2$ 上的作用. 此时商 $\mathbb A^2\sslash\mathbb G_m\cong\mathbb A^1$ 光滑, 但不是几何商. 按上述过程, 在原点爆破得到 $\operatorname{Bl}_0\mathbb A^2$, 并除去 $V(xy)$ 的严格变换. $\mathbb G_m$ 在补集 $U_1$ 上的作用稳定子有限, $U_1\to\mathbb A^2$ 是 $\mathbb G_m$-不变双有理映射 (但既不固有也不满射), 在商上诱导同构 $U_1/\mathbb G_m\to\mathbb A^2\sslash\mathbb G_m$.
\end{example}
